\title{Parameter-Efficient Orthogonal Finetuning via Factorization}

\begin{abstract}
Large-scale foundation models are becoming widespread, but training them from the ground up is prohibitively costly. Therefore, it is crucial to develop efficient methods to adapt these powerful models to downstream tasks. In this work, we investigate a systematic finetuning approach—Orthogonal Finetuning (OFT)—for adapting models to specific tasks. Although OFT achieves strong generalization, it still requires a relatively large number of trainable parameters because of the high dimensionality inherent in orthogonal matrices. To overcome this limitation, we analyze OFT through the lens of information transmission and identify several key criteria that can improve parameter efficiency. Drawing inspiration from the Cooley-Tukey fast Fourier transform algorithm, which facilitates efficient information transmission, we introduce a compact orthogonal parameterization based on butterfly structures. Incorporating this parameterization into OFT yields a novel, parameter-efficient finetuning technique, termed Orthogonal Butterfly (BOFT). BOFT generalizes OFT as a special instance, thereby offering a broader orthogonal finetuning framework. Finally, we perform comprehensive empirical evaluations by adapting large vision transformers, large language models, and text-to-image diffusion models to various downstream vision and language tasks.
\end{abstract}

\section{Introduction}

Recent advancements such as ChatGPT~\cite{brown2020language,chen2021evaluating} and Stable Diffusion~\cite{rombach2022high} highlight the extraordinary generalization capabilities of large foundation models. The rapid improvements in these models' performance have been accompanied by a significant increase in parameter counts (e.g., GPT-3 contains approximately 175 billion parameters). Consequently, training such foundation models from scratch has become increasingly difficult. Progress in the field thus hinges on the ability to adapt existing models efficiently without full retraining. Specifically, it is essential to develop approaches that allow effective adaptation of pre-trained foundation models to downstream tasks. Broadly, three primary strategies exist for efficient task adaptation: \emph{model finetuning}~\cite{hulora2022,zaken2022bitfit,guo2020parameter,raffel2020exploring,qiu2023controlling,zhang2022adaptive,chavan2023one}, which keeps the model architecture intact while finetuning a subset of its parameters; \emph{adapter tuning}~\cite{rebuffi2017learning,houlsby2019parameter,pfeiffer2020adapterfusion,lin2020exploring,he2021towards}, which introduces additional trainable parameters integrated into the original model; and \emph{prompt tuning}~\cite{li2021prefix,lester2021power}, which prepends trainable prefix tokens to the input. Among these, model finetuning stands out as a straightforward yet effective technique that, importantly, does not incur any increase in inference latency.

The core idea behind model finetuning is to maintain similarity between the pretrained model and the finetuned model according to specific metrics, thereby preserving the knowledge acquired during pretraining. For example, many existing finetuning techniques utilize a small learning rate as a heuristic to keep the deviation between the pretrained and finetuned models minimal. Considering the structured characteristics of weight matrices, a more rigorous strategy aims to retain the relational information within the weight matrices, specifically the pairwise angles between neurons. This concept gives rise to a novel finetuning framework called Orthogonal Finetuning~(OFT)~\cite{qiu2023controlling}, where neurons within the same layer undergo transformation via a shared orthogonal matrix, ensuring that the pairwise angles among neurons are theoretically maintained throughout the finetuning phase. Although OFT has shown encouraging generalization and convergence properties when applied to finetuning text-to-image diffusion models~\cite{qiu2023controlling}, the sheer dimensionality of orthogonal matrices results in a large number of trainable parameters. To mitigate this, OFT employs a block-diagonal structure to decrease the parameter count. Nevertheless, this parameter efficiency introduces drawbacks—the orthogonal matrix is constrained by a fixed sparsity pattern, and the orthogonal transformation is independently applied across separate blocks. While this block-diagonal sparsity pattern reduces parameters, it may induce undesirable inductive biases (e.g., the block-diagonal orthogonal matrix limits expressiveness and cannot replicate traditional linear transforms). Crucially, identifying an optimal sparsity pattern remains an unresolved challenge.

The central challenge in tackling this issue is to create a dense orthogonal matrix that remains parameter-efficient. Although at first this appears impractical because a $d$-dimensional dense orthogonal matrix traditionally requires $\mathcal{O}(d^2)$ parameters, we introduce a novel strategy that constructs a dense orthogonal matrix by combining multiple factorized sparse matrices. This method is motivated by the insight that the number of trainable parameters can be decreased by exchanging computational effort for memory savings. Since the orthogonal matrix is expressed as a product of sparse matrices, the multiplication must be repeated during each training step. To better understand the matrix factorization, we view the generation of a dense orthogonal matrix as an information transmission problem. Within this framework, producing a dense orthogonal matrix via multiplication of sparse matrices can be interpreted as the transmission of information across a grid-structured graph. This perspective allows us to devise numerous approaches for sparse matrix factorization that restrict the count of trainable parameters while maintaining sufficient expressiveness to produce dense matrices.

To attain parameter efficiency in our information transmission framework, we draw inspiration from the butterfly patterns found in the Cooley-Tukey fast Fourier transform algorithm~\cite{cooley1965algorithm}, where information exchange between nodes is highly efficient~\cite{klasing1994broadcasting}. In particular, the butterfly graph used in the Cooley-Tukey algorithm inherently defines a method for sparse matrix factorization known as butterfly factorization. Utilizing butterfly factorization, it becomes possible to generate a $d$-dimensional dense matrix as a product of $\mathcal{O}(\log d)$ sparse matrices, each containing $\mathcal{O}(d)$ non-zero elements. By ensuring the orthogonality of each sparse matrix, we obtain an efficient orthogonal parameterization with only $\mathcal{O}(d\log d)$ parameters, representing a substantial reduction compared to the original parameter count. Building on this efficient orthogonal parameterization, we introduce a new parameter-efficient finetuning technique called Orthogonal Butterfly~(BOFT). BOFT generalizes OFT as a special case and establishes a comprehensive orthogonal finetuning framework. A common feature of both the block-diagonal and butterfly structures is sparsity. Each introduces data sparsity into orthogonal matrices to diminish the effective number of trainable parameters. It is also insightful to contrast our method with the low-rank structure utilized in LoRA~\cite{hulora2022}; both low-rank and sparse matrices represent major families of structured matrices~\cite{candes2011robust} that enable parameter efficiency.

In contrast to the block-diagonal structure employed by OFT to balance expressiveness and regularity, BOFT utilizes the butterfly structure to create a more gradual interpolation between matrices from the full orthogonal group (expressivity) and identity matrices (regularity). This approach allows us to identify a smaller hypothesis class within the orthogonal group suitable for downstream tasks. Given the prevalent application of the butterfly structure in numerous fast linear transforms, including the discrete Fourier and discrete cosine transforms, we contend that our structured parameter-efficient approach introduces an implicit inductive bias that enhances generalizability and mitigates overfitting. Our reasoning is based on the fact that the butterfly structure can readily reconstruct many classical linear transforms, a capability that the block-diagonal structure in OFT lacks. Our contributions are summarized as follows:

\begin{itemize}[leftmargin=*,nosep]
    \item We investigate parameter efficiency in orthogonal finetuning through a novel information transmission framework. This framework recasts the challenge of designing a parameter-efficient dense orthogonal matrix as an information transmission problem on a grid-structured graph.
    \item Drawing inspiration from the butterfly structures found in the Cooley-Tukey algorithm, we introduce Orthogonal Butterfly, a parameter-efficient orthogonal finetuning approach, within the proposed information transmission framework.
    \item We offer several theoretical insights explaining why BOFT can substantially reduce the number of trainable parameters while maintaining robust expressivity and training stability. Additionally, due to the matrix factorization, BOFT supports an interesting weight interpolation (refer to Figure~\ref{fig:boft_interpolation}).
    \item For the first time, we apply orthogonal finetuning across a variety of tasks beyond controllable text-to-image generation~\cite{qiu2023controlling}, demonstrating its strong potential as a universal model finetuning technique. Specifically, we implement BOFT in diverse adaptation scenarios spanning computer vision and natural language processing. BOFT surpasses current state-of-the-art methods by a significant margin, confirming its superior parameter efficiency and generalization performance.
\end{itemize}

\section{Related Work}

\textbf{Parameter-efficient finetuning (PEFT)}. As foundational models (\emph{e.g.}, \cite{brown2020language,radford2021learning,rombach2022high,kirillov2023segment}) continue to grow in size and capability, there has been a surge of interest in finetuning these models for downstream applications in a parameter-efficient manner~\cite{treviso2023efficient,ding2023parameter,lialin2023scaling}. Among the various PEFT techniques~\cite{houlsby2019parameter,aghajanyan2020intrinsic,hulora2022,edalati2022krona,wang2022adamix,gheini2021cross,zaken2022bitfit,guo2020parameter,sung2021training,ansell2022composable,lester2021power,li2021prefix,vu2022spot,he2021towards,mao2021unipelt,karimi2021compacter,liu2022few,sung2022lst,chen2022parameter,jia2022visual,chen2022adaptformer,zhang2022neural,jie2023fact,lian2022scaling,luo2023towards}, reparameterization-based approaches~\cite{aghajanyan2020intrinsic,hulora2022,edalati2022krona} are most closely related to our study. LoRA~\cite{hulora2022} enhances pretrained weight matrices by adding a product of two low-rank matrices, showing promising results on natural language benchmarks. Because LoRA fixes the rank of these matrices across layers, several methods \cite{zhang2022adaptive,valipour2022dylora,zhang2023increlora} adaptively vary the rank per layer to optimize parameter allocation. Owing to its simplicity, this low-rank weight reparameterization technique has become widely adopted \cite{dettmers2023qlora,zi2023delta,chavan2023one}. Drawing inspiration from how hyperspherical energy relates to generalization~\cite{liu2018learning,liu2021orthogonal}, \cite{qiu2023controlling} introduced orthogonal finetuning (OFT) as an alternative, effective method for finetuning text-to-image diffusion models. OFT specifically learns an orthogonal matrix that transforms neurons within the same layer, yielding better generalization and more stable training compared to LoRA. Despite its advantages, OFT typically involves a larger number of trainable parameters than LoRA. Therefore, improving OFT’s parameter efficiency is a valuable objective. Additionally, it remains unclear if OFT can be effectively applied beyond controlling text-to-image diffusion models~\cite{qiu2023controlling}. BOFT enhances OFT’s parameter efficiency by employing butterfly factorization. This advancement enables us to showcase the effectiveness of orthogonal finetuning across a broader range of adaptation tasks.

\textbf{Butterfly structures}. The radix-2 Cooley-Tukey algorithm recursively decomposes an $N$-point discrete Fourier transform into two $\frac{N}{2}$-point discrete Fourier transforms, producing a butterfly structure expressible as a product of multiple sparse matrices (commonly referred to as a butterfly matrix). Butterfly matrices have been employed to parameterize orthogonal matrices to eliminate the need for pivoting in Gaussian elimination and enhance efficiency~\cite{parker1995random}, to stabilize the training of recurrent neural networks~\cite{jing2017tunable}, and in kernel approximation~\cite{munkhoeva2018quadrature}. The works \cite{dao2019learning,dao2019kaleidoscope} learn fast linear transforms via butterfly parameterizations. Additionally, \cite{chen2022pixelated,dao2022monarch} leverage butterfly matrices to facilitate efficient neural network training. Butterfly structures have also proven beneficial in fast matrix-vector multiplication~\cite{michielssen1996multilevel,o2010algorithm}, data-sparse matrix approximations~\cite{li2015butterfly}, and network communication~\cite{klasing1994broadcasting,li2003linear,rajasingh2016transmission}. Differing from these prior studies, our focus is on exploiting butterfly structures to improve the parameter efficiency of OFT.

\section{An Information Transmission View on Orthogonal Finetuning}

We begin by reviewing some fundamentals of OFT. To fine-tune the pretrained weight matrix, OFT reparameterizes the updated weight matrix as the product of a learnable orthogonal matrix and the fixed pretrained weight matrix. Unlike LoRA, which updates weights by adding a low-rank matrix, OFT modifies the weights multiplicatively using an orthogonal matrix. For parameter-efficiency, LoRA employs a low-rank structure, whereas the original OFT adopts a block-diagonal orthogonal structure~\cite{qiu2023controlling} (the smaller each diagonal block is, the more parameter-efficient OFT becomes). An intuitive comparison is shown in Figure~\ref{fig:comp_lora}. The rationale for applying orthogonal transformations to fine-tune the weight matrix is to maintain the pairwise angles among neurons~\cite{liu2018learning,liu2021orthogonal,qiu2023controlling}, thus largely preserving the semantic knowledge acquired during pretraining. Specifically, OFT optimizes an orthogonal matrix $\bm{R}\in\mathbb{R}^{d\times d}$ for a pretrained linear layer $\bm{W}^0\in\mathbb{R}^{d\times n}$, modifying the forward pass from $\bm{z}=(\bm{W}^0)^\top\bm{x}$ to $\bm{z}=(\bm{R}\bm{W}^0)^\top\bm{x}$, where $\bm{x}\in\mathbb{R}^d$ and $\bm{z}\in\mathbb{R}^n$ denote the input and output vectors, respectively. To enable zero initialization, OFT starts with $\bm{R}$ as the identity matrix. To guarantee $\bm{R}$ remains orthogonal during fine-tuning, we follow \cite{qiu2023controlling,liu2021orthogonal} and apply the Cayley parameterization, i.e., $\bm{R} = (\bm{I} + \bm{Q})(\bm{I} - \bm{Q})^{-1}$, where $\bm{Q}$ is a skew-symmetric matrix satisfying $\bm{Q} = -\bm{Q}^\top$. For parameter-efficiency, the block-diagonal form is introduced by expressing $\bm{R}$ as $\text{diag}(\bm{R}_1, \bm{R}_2, \ldots, \bm{R}_r)$, where each $\bm{R}_i \in \mathbb{R}^{b \times b}$ is a small orthogonal matrix and $br = d$. The parameter-efficiency gained from the block-diagonal structure comes at a cost: it imposes the assumption that the dimensions of a neuron (i.e., a column vector of $\bm{W}^0$) are partitioned into $r$ groups, and dimensions from different groups are transformed independently by distinct orthogonal matrices. Although the block-diagonal structure has been empirically successful, dividing neuron dimensions into $r$ groups purely based on their indices is not conceptually meaningful, which motivates the use of dense orthogonal matrices. This leads to a natural question: \emph{Is it possible to construct a dense orthogonal matrix while maintaining parameter-efficiency?}

To tackle this problem, we suggest constructing a dense orthogonal matrix as a product of several sparse orthogonal matrices. To offer a cohesive and intuitive framework for analyzing the sparsity structure in orthogonal matrix factorization, we reinterpret the task of generating a dense orthogonal matrix in OFT as an information transmission problem. More precisely, generating a dense matrix $\bm{R}\in\mathbb{R}^{d\times d}$ through a product of $m$ square matrices \(\thickmuskip=2mu \medmuskip=2mu \bm{R}=\bm{B}_m\bm{B}_{m-1}\cdots\bm{B}_1\) can be viewed as transmitting information across a grid composed of \(d \times (m+1)\) nodes, as depicted in Figure~\ref{fig:info_trans}. This information transmission perspective is inspired by the observation that a \(d\)-dimensional dense square matrix can be interpreted as representing dense connectivity between one set of \(d\) nodes and another. For the matrix \(\bm{R}\), if the entry \(R_{ij}\) equals zero, this implies that information from the \(j\)-th node cannot flow to the \(i\)-th node. Conversely, if \(R_{ij}\) is non-zero, information can be conveyed. Hence, expressing the dense matrix \(\bm{R}\) as a product of matrices \(\bm{B}_m\bm{B}_{m-1}\cdots\bm{B}_1\) can also be seen as sequentially exchanging information guided by the graphs defined by each \(\bm{B}_i, \forall i\). The flow of information occurs first along \(\bm{B}_1\) and finally through \(\bm{B}_m\). As a specific example illustrated in Figure~\ref{fig:info_trans}, we consider the factorization \(\bm{R}=\bm{B}_5\bm{B}_4\bm{B}_3\bm{B}_2\bm{B}_1\), whose sparsity patterns and corresponding induced graph are displayed. The graph in Figure~\ref{fig:info_trans} results from unfolding the matrix multiplication. Within this induced graph, each matrix \(\bm{B}_i\) is interpreted as the connectivity matrix from the nodes at level \(i\) to those at level \(i+1\). More specifically, the element \((j_1,j_2)\) of \(\bm{B}_i\) indicates whether there is a directed edge from the \(j_2\)-th node at level \(i\) to the \(j_1\)-th node at level \(i+1\) (with zero signifying no edge). For \(\bm{B}_5\bm{B}_4\bm{B}_3\bm{B}_2\bm{B}_1\) to be a dense matrix, every node at the first level must be capable of transmitting information to all nodes at the sixth level. Restricting attention to \(\bm{R}=\bm{B}_4\bm{B}_3\bm{B}_2\bm{B}_1\), which corresponds to source nodes at level one and receiver nodes at level five, we observe that information from node 1 cannot reach node 3. Consequently, the sparsity pattern of \(\bm{B}_4\bm{B}_3\bm{B}_2\bm{B}_1\) contains a zero at position \((3,1)\). A similar relationship between the induced graph and the sparsity pattern holds for \(\bm{R}=\bm{B}_3\bm{B}_2\bm{B}_1\) and \(\bm{R}=\bm{B}_2\bm{B}_1\). More generally, for a matrix \(\bm{R}\in\mathbb{R}^{d\times d}\) to be dense, the \(m\) factor matrices \(\bm{B}_m,\ldots,\bm{B}_1\) must correspond to a set of directed edges arranged on a \(d \times (m+1)\) grid, where each directed edge connects nodes only between adjacent levels (i.e., columns), ensuring that information from every node at the first level can reach every node at the \((m+1)\)-th level.

The perspective of information transmission aids in deepening our understanding of the sparsity pattern found in factorization matrices within OFT. Figure~\ref{fig:blk_diag} displays the block-diagonal arrangement of $\bm{R}$ in the original OFT. Although this block-diagonal formation decreases the number of trainable parameters, it is incapable of creating a dense matrix $\bm{R}$. Our objective is to form a dense orthogonal matrix by multiplying $m$ sparse orthogonal matrices, while minimizing the number of effective trainable parameters. From the information transmission standpoint, the primary requirements to achieve this are \emph{(i)} \textbf{dense connectivity}: every node in the first layer must have at least one connecting path to every node in the final layer, and \emph{(ii)} \textbf{minimum free edges}: the total number of edges should be as low as possible given the orthogonality constraint. It is important to recognize that orthogonality imposes a subtle constraint on the edges between adjacent layers. For instance, for each matrix $\bm{B}_i$ to be full-rank (a prerequisite for orthogonality), there must be at least $d$ edges establishing a one-to-one correspondence between all nodes in the $i$-th layer and those in the $(i+1)$-th layer, making the edge count between these layers at least $d$ (for example, 4 in Figure~\ref{fig:info_trans}). These $d$ edges are essential for orthogonality and should be excluded from the total edge count because these elements are not trainable (e.g., in a $d\times d$ orthogonal matrix with $d$ non-zero entries, these entries are restricted to $\pm 1$). Since orthogonal matrices require fewer parameters than full matrices, the orthogonality constraint introduces additional dependencies among edges. As an illustration, in a $2\times 2$ orthogonal matrix, having a zero at position $(1,1)$ implies a zero at position $(2,2)$ (i.e., one missing edge can indicate another missing edge). Consequently, for each permissible edge configuration, orthogonality may sometimes necessitate adding or removing edges. Visualizing the non-zero pattern of the composed orthogonal matrix is especially beneficial in OFT because the focus is solely on the non-zero trainable components of $\bm{R}$, and the exact values of these components are not significant.

A straightforward dense connection between two levels requires $\mathcal{O}(d^2)$ edges (i.e., a single dense orthogonal matrix), resulting in $d^2 - d$ edges (for $d=4$, this equals 12 edges). Figure~\ref{fig:info_trans} provides an example of a viable matrix factorization, which uses a total of $10$ edges, actually fewer than a single dense orthogonal matrix. This approach allows us to analyze the parameter-efficiency of OFT through a graphical lens, enabling easy derivation of feasible factorizations. We draw motivation from an intriguing topology derived from the Cooley-Tukey algorithm, known as butterfly graphs~\cite{cooley1965algorithm}, which can efficiently create dense connections between $d$ source nodes and $d$ receiver nodes using only $\mathcal{O}(d\log d)$ edges. For instance, while the topology in Figure~\ref{fig:info_trans} requires $10$ edges for dense connectivity, the butterfly network only needs $8$ edges. In the following, we explain how the butterfly structure enhances parameter efficiency.

\section{Orthogonal Parameterization by Butterfly Factorization}

Originally developed for the Cooley-Tukey algorithm to execute the fast Fourier transform, the butterfly structure enables a local alteration in the frequency domain to induce a global change in the spatial domain. This concept is analogous to our information transmission scenario, where each node at the first level can send information to all nodes at the last level. The butterfly structure has also gained popularity as a computer network topology~\cite{solihin2015fundamentals,li2011linear} for efficient information exchange. Assuming $k \geq 2$ is a power of $2$, we first define the \emph{butterfly factor} $\bm{B}^F(k)$ as

\begin{equation}\label{eq:but_factor}
\begin{aligned}
    \bm{B}^F(k)=\begin{bmatrix}
    \text{diag}(\bm{d}_1) & \text{diag}(\bm{d}_2) \\
    \text{diag}(\bm{d}_3) & \text{diag}(\bm{d}_4)
    \end{bmatrix}\in\mathbb{R}^{k\times k},
\end{aligned}
\end{equation}

where each $\bm{d}_i \in \mathbb{R}^{\frac{k}{2}}$ for all $i$ are vectors. Given $d=2^N$, the $d$-dimensional \emph{butterfly matrix} $\bm{B}(d) \in \mathbb{R}^{d \times d}$ is then defined recursively as

\begin{equation}
\begin{aligned}
    \bm{B}(d) = \tilde{\bm{B}}(d,d) \cdot \begin{bmatrix}
    \bm{B}_1(\frac{d}{2}) & \bm{0} \\
    \bm{0} & \bm{B}_2(\frac{d}{2})
    \end{bmatrix} = \tilde{\bm{B}}(d,d) \tilde{\bm{B}}(d,\frac{d}{2}) \cdots \tilde{\bm{B}}(d,2),
\end{aligned}
\end{equation}

where $\bm{B}_1(\frac{d}{2})$ and $\bm{B}_2(\frac{d}{2})$ represent two butterfly matrices each of dimension $\frac{d}{2}$. We then introduce the \emph{butterfly component} defined as $\thickmuskip=2mu \medmuskip=2mu \tilde{\bm{B}}(d,k)=\text{diag}(\bm{B}^F_1(k),\cdots,\bm{B}^F_{d/k}(k))$, which is a block-diagonal matrix of size $d \times d$ with blocks of size $k$, where each $\bm{B}^F_i(k), \forall i$ corresponds to butterfly factors given in Equation~\ref{eq:but_factor}. At this point, we are prepared to utilize the butterfly matrix for parameterizing an orthogonal matrix. To do so, it suffices to ensure that all multiplicative factors $\tilde{\bm{B}}(d,k), \forall k$ composing the butterfly matrix $\bm{B}(d)$ are orthogonal. We first examine the block-diagonal matrix $\tilde{\bm{B}}(d,2)$ with block size $2$, and it is straightforward to guarantee the orthogonality of $\tilde{\bm{B}}(d,2)$ by applying the Cayley transform (or a $2$-dimensional rotation) for parameterizing each block, following approaches in \cite{liu2021orthogonal,qiu2023controlling}. The pattern of non-zero entries in every butterfly component corresponds to a permutation of the non-zero pattern of $\tilde{\bm{B}}(d,2)$; hence, all butterfly components can be naturally parameterized as orthogonal matrices. This provides an efficient scheme for orthogonal matrix parameterization constructed from numerous $2 \times 2$ orthogonal matrices. Extending the butterfly matrices as in \cite{chen2022pixelated}, we define a block butterfly component $\tilde{\bm{B}}^b(d,k)$ where each entry in $\bm{d}_i, \forall i$ is replaced by a $b \times b$ matrix. To ensure orthogonality of the block butterfly component $\tilde{\bm{B}}^b(d,2)$, we parameterize each $2b \times 2b$ block matrix to be orthogonal. The non-zero structures of other butterfly components $\tilde{\bm{B}}^b(d,k)$ for $k > 2$ correspond to block-wise permutations of that of $\tilde{\bm{B}}^b(d,2)$, and thus can also be similarly parameterized as orthogonal matrices. Putting these components together, the forward pass in BOFT is given by

\begin{equation}
    \begin{aligned}\nonumber
    \bm{z}=\big(\bm{R}(m,b)\cdot\bm{W}^0\big)^\top\bm{x},~~~~\text{s.t.}~\bigg{\{}\bm{R}(m,b)=\prod_{i=1}^m \tilde{\bm{B}}^b_{(d,i)}~~\&~~\underbrace{\big(\tilde{\bm{B}}^b_{(d,j)}\big)^\top\tilde{\bm{B}}^b_{(d,j)}=\tilde{\bm{B}}^b_{(d,j)}\big(\tilde{\bm{B}}^b_{(d,j)}\big)^\top\!=\bm{I}_d}_{\forall j\in [1, m]}\bigg{\}},
    \end{aligned}
\end{equation}

Here, for simplicity, we denote $\tilde{\bm{B}}^b(d,2^{m - i+1})$ as $\tilde{\bm{B}}^b_{(d,i)}$, and $\bm{I}_d$ represents the $d$-dimensional identity matrix. The orthogonal matrix $\bm{R}(m,b)\in\mathbb{R}^{d\times d}$ is formed by multiplying several orthogonal butterfly components. For ease of notation, BOFT with $\bm{R}(m,\frac{b}{2})$ is referred to as BOFT($m$,$b$), where $b\geq 2$. When $m=1$, BOFT($1$,$b$) corresponds to the block-diagonal OFT~\cite{qiu2023controlling} with block size $b$. BOFT($1$,$d$) simplifies to the original OFT~\cite{qiu2023controlling} having a fully unconstrained orthogonal matrix. The setting BOFT($\log \frac{2d}{b}$,$b$) employs the block butterfly matrix $\bm{B}^{\frac{b}{2}}(d)$ as $\bm{R}$, producing a dense orthogonal matrix $\bm{R}$. Generally, BOFT($m$,$b$) involves $\frac{1}{2}(b-1)dm$ effective trainable parameters for fine-tuning a linear layer of dimension $d\times n$. Using the butterfly matrix, i.e., $m=\log d, b=2$, BOFT requires $\mathcal{O}(d\log d)$ parameters. By comparison, the original OFT with a fully dense orthogonal matrix demands $\mathcal{O}(d^2)$ parameters, while the block-diagonal OFT with $r$ blocks needs $\mathcal{O}(bd)$ parameters. Consequently, the original OFT must select block size $b=d$ to generate a dense orthogonal matrix, whereas BOFT can employ any $b$ to achieve this goal.

\textbf{Identity initialization for BOFT}. Fine-tuning approaches typically begin from the pretrained model itself to ensure the fine-tuned version remains close to the pretrained parameters. For instance, LoRA initializes the low-rank weights at zero. In BOFT, all butterfly components are initialized as identity matrices (equivalently, the skew-symmetric matrices used in the Cayley transform are initialized to zero).

\textbf{Multiplicative Dropout for BOFT}. LoRA~\cite{hulora2022} additionally introduces a Dropout layer on the low-rank weight updates to mitigate overfitting. Standard Dropout~\cite{srivastava2014dropout} functions naturally for LoRA, but is incompatible with BOFT due to its multiplicative weight updates. To overcome this, we propose a multiplicative Dropout method tailored for BOFT. Since the orthogonal matrix $\bm{R}(m,b)$ is formed by $m$ orthogonal butterfly components, each can be rearranged into $2b\times 2b$ block-diagonal orthogonal matrices. This multiplicative Dropout randomly selects $p_1$ fraction of the butterfly components and $p_2$ fraction of the diagonal blocks within each butterfly component, replacing these blocks with identity matrices.

\section{Intriguing Insights and Discussions}

\textbf{Expressivity of BOFT}. The butterfly architecture combined with permutations can exactly reconstruct many traditional fast linear transforms~\cite{dao2019learning,dao2019kaleidoscope} (such as the fast Fourier transform and Hadamard transform). However, the capability of our orthogonal butterfly matrix to approximate a general orthogonal matrix is still unclear. We begin by running a simulation to approximate a random dense orthogonal matrix~\cite{anderson1987generation} of size $512 \times 512$. The outcomes shown in Figure~\ref{fig:para_eff} are averaged over 10 different random seeds. The vertical axis represents the approximation error, while the horizontal axis shows the number of effective trainable parameters. Each curve of the same color corresponds to BOFT with a fixed block size, where the leftmost point indicates the error of BOFT($1$, $b$) (i.e., the original block-diagonal OFT with block size $b$). Overall, BOFT tends to provide better parameter efficiency compared to OFT. For instance, BOFT($9$, $2$) exhibits greater expressiveness than BOFT($1$, $16$) while using significantly fewer parameters. Generally, BOFT with smaller $b$ and larger $m$ is more parameter-efficient; for example, BOFT($6$, $4$) utilizes considerably fewer parameters but achieves a similar approximation error to BOFT($2$, $16$). In summary, the butterfly matrix constitutes a more structured subset within the orthogonal group (relative to the block-diagonal form), making BOFT demonstrably more expressive than OFT for the same block size.

\begin{theorem}[Expressivity of BOFT]\label{thm:exp_boft}
    BOFT possesses greater expressivity than OFT when using the same block size. To approximate any orthogonal matrix of dimension $d$, one can express it as a product of butterfly matrices in the form $\bm{B}_{d-1,1}(d)\bm{B}^\top_{d-1,2}(d)\cdots\bm{B}_{1,1}(d)\bm{B}^\top_{1,2}(d)$, where $\bm{B}_{i,j}(d)$ are butterfly matrices for all $i,j$.
\end{theorem}

Theorem~\ref{thm:exp_boft} motivates a straightforward extension for BOFT—where the resulting orthogonal matrix is generalized as $\thickmuskip=2mu \medmuskip=2mu \bm{R}^G(m_1,b_1,m_2,b_2,l) = \bm{R}_{l,1}(m_1,b_1) \bm{R}_{l,2}^T(m_2,b_2) \cdots \bm{R}_{1,1}(m_1,b_1) \bm{R}_{1,2}^T(m_2,b_2)$, with $\bm{R}_{i,j}^T(m,b)$ denoting the orthogonal matrices employed in BOFT. If $m_1 = m_2 = \log d$, $b_1 = b_2 = 2$, and $l = d-1$, then $\bm{R}^G(m_1,b_1,m_2,b_2,l)$ can represent the full orthogonal group. This type of matrix composition is also known as the kaleidoscope hierarchy~\cite{dao2019kaleidoscope}. Nonetheless, it is important to note that increased expressiveness does not necessarily translate into improved finetuning performance, since full finetuning, despite its universal expressiveness, often produces unsatisfactory results. Balancing expressivity with regularization is crucial for the generalizability of model finetuning. BOFT expands the finetuning parameter space while incorporating structural priors, allowing for a better balance between expressivity and regularization.

\textbf{Spectral properties}. Orthogonal finetuning typically produces superior spectral characteristics compared to LoRA, as it exactly maintains the spectral norm of the original pretrained weight matrix $\bm{W}^0$. This can be understood through singular value decomposition: $\bm{W}^0=\bm{U}\bm{\Sigma}\bm{V}^\top$, where $\bm{U}$ and $\bm{V}$ are orthogonal matrices and $\bm{\Sigma}$ is a diagonal matrix of singular values. Both OFT and BOFT apply an orthogonal matrix $\bm{R}$ on the left side, resulting in the finetuned weights $\bm{R}\bm{U}\bm{\Sigma}\bm{V}^\top$, which preserves the largest singular value (i.e., the spectral norm of $\bm{W}^0$). This preservation has been demonstrated to significantly improve training stability and generalization~\cite{miyato2018spectral,yoshida2017spectral}. Additional intriguing mathematical properties are discussed in Appendix~\ref{app:sn_property}.

\textbf{Orthogonal finetuning as learning bilinear similarity}. BOFT can be expressed as learning the bilinear similarity $\bm{w}^0_i\bm{R}\bm{x}$, where $\bm{w}^0_i$ represents the $i$-th neuron (i.e., column vector) of the weight matrix $\bm{W}^0$. In this perspective, BOFT amounts to learning the bilinear similarity matrix $\bm{R}$ under a strong constraint (i.e., $\bm{R}$ must be orthogonal), which inherently relates to distance metric learning~\cite{xing2002distance} and bilinear forms~\cite{roman2005advanced}.

\textbf{Inductive bias and generalization in BOFT}. Since $\bm{R}(m,b)$ in BOFT typically corresponds to a structured subset of the orthogonal group that limits the hypothesis space, BOFT inherently introduces an inductive bias. We contend that the structured inductive bias induced by the butterfly factorization enhances generalization, as it embodies a common structured pattern found in numerous classical linear transforms~\cite{dao2019kaleidoscope}, such as the discrete Fourier transform, discrete sine/cosine transforms, and Hadamard transform. Furthermore, the sparse matrix factorization in BOFT might also contribute implicit inductive biases~\cite{gunasekar2017implicit,li2018algorithmic,liu2019neural}.

\textbf{Comparison to butterfly-based sparse training}. Several studies~\cite{dao2019learning,dao2019kaleidoscope,chen2022pixelated,dao2022monarch} have explored sparse training utilizing the butterfly parameterization. These works generally emphasize directly reparameterizing weight matrices through the butterfly framework and training neural networks from the beginning. In contrast, \cite{dao2022monarch} investigates finetuning pretrained weights by initially projecting them onto a variant of butterfly matrices, followed by optimizing these projected elements for downstream tasks. BOFT introduces a distinctly different finetuning approach, which modifies the weights using layer-shared weight matrices.

\input{rebuttal-results}

\section{Concluding Remarks and Limitations}

This paper presents Orthogonal Butterfly, a versatile and parameter-efficient finetuning technique for foundation models grounded in the butterfly structure. The fundamental idea for achieving enhanced parameter efficiency lies in representing a dense orthogonal matrix as the product of several sparse orthogonal matrices. To facilitate the discovery of viable matrix factorizations, we introduce a graphical information transmission framework. Within this framework, we identify that the butterfly structure effectively satisfies our criteria for sparse orthogonal matrix factorization. We validate the practical efficacy of BOFT in finetuning large language models, expansive vision models, and text-to-image generative models. Our experimental results further confirm BOFT’s superiority as a broadly applicable finetuning method.

Despite its demonstrated empirical success, BOFT is not without limitations. Because the final orthogonal matrix in BOFT is composed as a product of multiple orthogonal matrices, its training runtime incurs a slightly greater overhead compared to OFT. Enhancing BOFT’s training efficiency remains an open challenge. Fortunately, once finetuning is complete, the orthogonal matrices learned by BOFT can be directly incorporated into the pretrained model without introducing extra inference latency. Additionally, it remains unclear whether the butterfly network represents the most efficient architecture for information transmission. Our information transmission framework enables us to draw insights from a different field—computer networking—where the efficiency of network topologies for information flow is extensively studied. We anticipate that alternative network structures could yield more efficient compositions of dense orthogonal matrices.

\end{document}